% ======================================================================
% CONSENSUS SKELETON: one-column workshop/arXiv ML-systems paper
% Synthesized from three independent cluster analyses (10 papers):
%   [A] ML-methods cluster: 2511.10676 (MLSys), 2605.27081 (ICML),
%       2607.08780 StickyMoE (plain article, one column)
%   [B] systems cluster: Fate (acmart sigplan), LayerScope (acmart),
%       Pre-gated MoE (IEEEtran/ISCA)
%   [C] landmark cluster: OLMoE (NeurIPS), MoE-Infinity (ICML),
%       Switch Transformers (JMLR), MoE-Beyond (plain article)
%
% Target style (mandated): plain article, ONE column, 10pt, letterpaper,
% natbib NUMBERED citations, booktabs tables. Compiles with pdflatex on
% standard TeX Live; no venue .sty required.
% ======================================================================

% [all 3 clusters] every paper roots at \documentclass{article}.
% One column per target style (matches StickyMoE/OLMoE/Switch). For a
% two-column venue add [twocolumn] or swap in the venue style.
\documentclass[10pt,letterpaper]{article}

% ---- Geometry ---------------------------------------------------------
% [B] systems template uses geometry to approximate a conference block;
% StickyMoE [A] uses margin=1in. letterpaper matches the target.
\usepackage[letterpaper,margin=1in]{geometry}

% ---- Core preamble [MANDATORY: agreed by all 3 clusters] --------------
% Order follows cluster A consensus: microtype, graphicx, subcaption,
% booktabs, amsmath ... hyperref last.
\usepackage[T1]{fontenc}
\usepackage[utf8]{inputenc}
\usepackage{microtype}          % all 3 clusters
\usepackage{graphicx}           % all 3 clusters
\usepackage{subcaption}         % all 3 clusters (A,C) / subfig in B-P3
\usepackage{booktabs}           % MANDATORY: booktabs-only tables, no \hline
\usepackage{amsmath}            % all 3 clusters
\usepackage{amssymb}            % A-P2/P3, B-P3, C-OLMoE
\usepackage{mathtools}          % A-P2/P3, C-OLMoE
\usepackage{amsthm}             % A-P2/P3 (theorem suite)
\usepackage{xcolor}
\usepackage{xspace}             % B-P3, C (system-name macros)
\usepackage{enumitem}           % B-P2, C (compact contribution lists)
\usepackage[hyphens]{url}       % B-P3, A-P2 (URL breaking)

% ---- hyperref: colorlinks with muted colors [MANDATORY] ----------------
% A: StickyMoE explicit colored setup. C: OLMoE single dark blue /
%    Switch all black. B: acmart screen colors / hyperref default.
% Consensus: colorlinks=true, muted (no default red boxes).
\definecolor{linkcolor}{RGB}{0,0,128} % C-OLMoE dark blue
\usepackage{hyperref}
\hypersetup{
  colorlinks = true,
  linkcolor  = blue!70!black,   % A-StickyMoE
  citecolor  = green!50!black,  % A-StickyMoE (alt: single linkcolor, C-OLMoE)
  urlcolor   = linkcolor,
}

% ---- natbib NUMBERED [target style] ------------------------------------
% All 3 clusters agree on BibTeX (never biblatex) with semantic keys.
% C-OLMoE passes numbers,compress to natbib; numbered per target style.
% Alternatives: A-P1/P2 venue-provided natbib; B numeric via cite/acmart;
% C-Switch author-year natbib (JMLR). Use \citep/\citet (A-P3, C).
\usepackage[numbers,sort&compress]{natbib}

% ---- cleveref (optional; A-P2/P3, C-OLMoE/MoE-Infinity) ----------------
\usepackage[capitalize,noabbrev]{cleveref}

% ---- Theorem environments [A-P2 full suite, A-P3 subset] ---------------
\theoremstyle{plain}
\newtheorem{theorem}{Theorem}[section]
\newtheorem{proposition}[theorem]{Proposition}
\newtheorem{lemma}[theorem]{Lemma}
\theoremstyle{definition}
\newtheorem{definition}[theorem]{Definition}
\theoremstyle{remark}
\newtheorem{remark}[theorem]{Remark}

% ---- Notation macros ----------------------------------------------------
% [A-P3, C] \textsc system-name macro with \xspace. Bold-vector math is
% A-P3-only; plain italic is the A/C majority -- keep macros minimal.
\newcommand{\method}{\textsc{<<METHODNAME>>}\xspace}
\newcommand{\R}{\mathbb{R}}
\newcommand{\calL}{\mathcal{L}}
% hyperref+algorithm coexistence fix (A-P1/P2, C-MoE-Infinity):
\newcommand{\theHalgorithm}{\arabic{algorithm}}

% ======================================================================
\begin{document}

% ---- Title block: plain \maketitle (A-P3, B-fallback, C-plain) ---------
\title{\method: <<TITLE>>}
\author{<<AUTHOR>> \\ <<AFFILIATION>> \\ \texttt{<<EMAIL>>}}
\date{} % suppress date (B-P1, all templates)
\maketitle

% ---- Abstract [MANDATORY rhetoric, all 3 clusters] ---------------------
% ONE paragraph: context/problem -> prior-work gap -> "We propose X ..."
% + mechanism -> BOLDED headline numbers -> artifact URL.
\begin{abstract}
<<CONTEXT: MoE models ...; PROBLEM/GAP: existing approaches ....>>
We propose \method, <<ONE-LINE DESCRIPTION>>. <<KEY INSIGHT + 2-3
techniques in order>>. On <<MODELS/HARDWARE>>, \method achieves
\textbf{<<XX\% / X.X$\times$>>} improvement over <<SOTA BASELINE>>.
\noindent\textbf{Code}: \url{https://github.com/<<REPO>>}
\end{abstract}

% ----------------------------------------------------------------------
\section{Introduction}
\label{sec:intro}
% [all 3 clusters] context -> problem -> categorize prior remedies ->
% "we propose" paragraph naming components "First ... Second ... Third ..."
% (A-P1, B) or "(i) ... (ii) ..." (A-P2) -> results-summary paragraph.
<<INTRO PARAGRAPHS>>

% Motivational page-1 figure/table is common (B-P2, C-OLMoE/MoE-Infinity).
\begin{figure}[t] % [t]-family placement MANDATORY (all 3 clusters)
  \centering % \centering in floats MANDATORY
  \includegraphics[width=\linewidth]{figures/<<MOTIVATION_FIG>>}
  \caption{<<Caption BELOW figure (MANDATORY). Optional bold lead-in
    sentence (A-P2, C-OLMoE).>>}
  \label{fig:motivation} % fig: prefix MANDATORY; \label right after \caption
\end{figure}

% Intro ENDS with an explicit itemized contributions list [MANDATORY].
% Evidence: A-P3 (enumerate), all of B (itemize of bold-tagged bullets),
% C 3/4 (Switch itemize; MoE-Infinity/MoE-Beyond bold run-in paragraphs).
% Alternative: bold run-in "\mypar{Contribution N}" style (C).
\paragraph{Contributions.}
We make the following contributions:
\begin{itemize}[leftmargin=1.5em]
  \item \textbf{<<Technique.>>} <<Core mechanism (\cref{sec:method}).>>
  \item \textbf{<<System.>>} <<Implementation in \method (\cref{sec:experiments}).>>
  \item \textbf{<<Evaluation.>>} <<Headline numbers vs.\ named baselines
        (\cref{sec:experiments}).>>
\end{itemize}

% ----------------------------------------------------------------------
\section{Background}
\label{sec:background}
% [all 3 clusters] preliminaries before method; numbered equations.
% B-P2 merges Related Work here ("Background and Related Work");
% B-P3 folds Motivation into Background.
\subsection{<<MoE PRIMER / PRELIMINARIES AND NOTATION>>}
\begin{equation}
  y = \sum_{i=1}^{N} G(x)_i \, E_i(x)
  \label{eq:moe} % eq: prefix MANDATORY (A-P2/P3, C-OLMoE/Switch)
\end{equation}

% ----------------------------------------------------------------------
\section{Related Work}
\label{sec:related}
% POSITION: early here, per target skeleton. DISAGREEMENT NOTE: B-P1/P3
% and C-OLMoE/Switch place Related Work LATE (after Evaluation); B-P2 and
% C-MoE-Infinity merge it early. Move this section after Evaluation if
% following the late convention.
% ORGANIZATION [all 3 clusters]: thematic categories, each a subsection
% (A-P3) or bold run-in paragraph (A-P2, B-P1); end by contrasting \method.
\subsection{<<Category A, e.g., System-Level Approaches>>}
\subsection{<<Category B, e.g., Post-Hoc Adaptation>>}
\subsection{<<Category C, e.g., Training-Time Methods>>}

% ----------------------------------------------------------------------
\section{\method: <<ONE-LINE METHOD DESCRIPTION>>}
\label{sec:method}
% [all 3 clusters] Overview first (with architecture figure), then one
% subsection per named component; one NUMBERED equation per loss term,
% each labeled eq:... and referenced from text.
\subsection{Overview}
\label{sec:method_overview}
<<OVERVIEW: name the 2-3 components as (i) ... (ii) ... (iii) ...>>
\subsection{<<Training Objective / Core Mechanism>>}
\begin{equation}
  \label{eq:objective}
  \calL = <<LOSS TERMS>> % one equation per loss term (A-P2/P3 mandatory)
\end{equation}
% Multi-row formulations: use align with \label on each row (A-P1).

% ----------------------------------------------------------------------
\section{Experiments}
\label{sec:experiments}
% [all 3 clusters] order: SETUP first -> main results -> analysis ->
% ABLATIONS LAST. Baselines are NAMED prior systems; quote x/% gains.
\subsection{Experimental Setup}
\label{sec:setup}
<<SETUP: models, datasets, hardware, workloads, hyperparameters.>>
Baselines: <<BASELINE 1>>~\citep{<<semantic-key-1>>},
<<BASELINE 2>>~\citep{<<semantic-key-2>>}.
\subsection{<<Main Results>>}
\begin{table}[t] % [t]-family placement MANDATORY
  \centering
  % Caption ABOVE table: MANDATORY (A-P2/P3, all of B, C-OLMoE; target
  % style). Alternatives: below table (A-P1, C-MoE-Infinity/Switch).
  \caption{\textbf{<<Bold lead-in sentence.>>} <<Self-contained details:
    comparison, protocol, takeaway (A-P2, C-OLMoE/Switch).>>}
  \label{tab:main} % tab: prefix MANDATORY
  \begin{tabular}{lcc}
    \toprule    % booktabs ONLY: \toprule/\midrule/\bottomrule, no \hline,
    <<Method>> & <<Metric 1>> ($\uparrow$) & <<Metric 2>> ($\downarrow$) \\ % no vertical rules
    \midrule    % (arrows on metrics: C-Switch)
    <<Baseline>> & -- & -- \\
    \midrule    % \midrule group separators for cohorts (C-OLMoE)
    \method & \textbf{<<best>>} & \textbf{<<best>>} \\ % bold best cells (all of C)
    \bottomrule
  \end{tabular}
  % Wide tables: \footnotesize wrapper (C-OLMoE) or
  % \resizebox{\linewidth}{!}{...} (B-P2, C-MoE-Infinity).
\end{table}
\subsection{<<Systems Evaluation / Performance Analysis>>}
% [B, A-P1] end-to-end performance, then accuracy/quality, then
% sensitivity / component drill-downs.
\subsection{Ablation Studies}
\label{sec:ablation} % ablations LAST within Experiments (all 3 clusters)

% ----------------------------------------------------------------------
\section{Discussion and Limitations}
\label{sec:discussion}
% Boundary/Limitations section [target skeleton; A-P1/P3, B-P3 discussion,
% C-MoE-Beyond explicit Limitations section, C-OLMoE appendix topic].
\paragraph{Limitations.}
<<LIMITATIONS: when \method does NOT help; scope boundary.>>

% ----------------------------------------------------------------------
\section{Conclusion}
\label{sec:conclusion}
<<Restate problem, method, key insight, headline numbers, future work.>>

% ---- Acknowledgments: unnumbered, before bibliography [all of B, C] ----
\section*{Acknowledgments}
<<FUNDING>>

% ---- Bibliography BEFORE appendix [MANDATORY, all 3 clusters] ----------
% Target style: plainnat + references.bib. COMMENTED OUT: no .bib ships
% with this template; uncomment in the real paper.
% Alternatives: venue bst (mlsys2025/icml2026/ACM-Reference-Format/
% IEEEtranS/acl_natbib) or B-P1 \bibliographystyle{plain}.
% \bibliographystyle{plainnat}
% \bibliography{references}

% ---- Appendix AFTER the bibliography [MANDATORY: A-P2/P3, B-P3, C 3/4] -
\appendix
\section{<<Appendix: Reproducibility / Proofs / Additional Results>>}
\label{app:repro} % app: prefix (A-P2/P3, C-Switch)
<<Reference from main text via \cref{app:repro}.>>

\end{document}
